\chapter{Marco teórico}
\label{cap:marco-teorico}

El marco teórico reúne el conocimiento previo necesario para entender el TFG.
No debe ser una colección de resúmenes independientes, sino una explicación
ordenada de conceptos, normativa, tecnologías y trabajos relacionados que
conduzca de forma natural al problema planteado.

\section{Marco legal, normativo o técnico}
\label{sec:marco-normativo}

Incluye en esta sección las normas, reglamentos, estándares, instrucciones
técnicas o criterios institucionales que condicionen el trabajo. En proyectos de
ingeniería puede ser necesario citar normativa de seguridad, reglamentos
industriales, estándares de diseño, documentación de fabricantes o guías
metodológicas.

Cuando una norma sea central para el trabajo, no basta con nombrarla: explica
qué requisito afecta a tu diseño, cálculo, instalación, simulación o evaluación.
Si el TFG no depende de normativa específica, utiliza este apartado para
delimitar los criterios técnicos que se tomarán como referencia.

\section{Estado de la cuestión}
\label{sec:estado-cuestion}

El estado de la cuestión resume las aportaciones previas más relevantes. Debe
permitir entender qué se ha hecho antes, qué limitaciones permanecen abiertas y
qué hueco ocupa tu TFG.

Las referencias bibliográficas se citan con el comando
\texttt{\textbackslash cite\{clave\}}. Por ejemplo, esta plantilla incluye
referencias clásicas sobre \LaTeX{} y composición de documentos
técnicos~\cite{lamport1994latex,mittelbach2004latex}. Sustituye estas entradas por
artículos, libros, normas, manuales o documentación técnica realmente empleados
en tu trabajo.

Evita basar el estado de la cuestión en fuentes poco fiables o sin trazabilidad.
En una memoria académica es preferible utilizar artículos científicos, libros,
normativa oficial, documentación técnica de fabricantes o informes
institucionales.


\section{Citas y referencias bibliográficas}
\label{sec:citas-referencias}

Toda afirmación que proceda de una fuente externa debe ir acompañada de una
cita. Esto incluye datos, definiciones, resultados de otros autores, imágenes,
tablas, normativa y criterios técnicos que no sean de elaboración propia.

La bibliografía se gestiona en el archivo \texttt{bibliography.bib}. Cada
entrada tiene una clave, como \texttt{knuth1984texbook}, que se usa dentro de
la memoria. Sólo aparecerán en la bibliografía final las referencias citadas en
el texto.

Si tu TFG utiliza siglas, preséntalas la primera vez que aparezcan. Por ejemplo:
Trabajo Fin de Grado (TFG). A partir de ese punto puede utilizarse la sigla.

\section{Figuras, tablas y referencias cruzadas}
\label{sec:referencias-cruzadas}

Las figuras, tablas, ecuaciones, capítulos y secciones deben referenciarse con
\texttt{\textbackslash label} y \texttt{\textbackslash ref}. Este sistema evita
errores de numeración cuando se añaden o eliminan elementos.

Ejemplos de prefijos recomendados para las etiquetas:

\begin{itemize}
    \item \texttt{cap:} para capítulos.
    \item \texttt{sec:} para secciones y subsecciones.
    \item \texttt{fig:} para figuras.
    \item \texttt{tab:} para tablas.
    \item \texttt{eq:} para ecuaciones.
    \item \texttt{lst:} para listados de código.
    \item \texttt{alg:} para algoritmos.
\end{itemize}

Antes de entregar, revisa que no queden referencias sin resolver. Si aparecen
signos de interrogación en el PDF, normalmente es necesario compilar varias
veces o corregir alguna etiqueta.
